% ============================================================================
% EXEMPLO FICTÍCIO — Relatório de TAL como Prática Profissional
% Curso Técnico Integrado em Informática (PPC 2012) — IFRN Campus Ceará-Mirim
% ============================================================================
%
% ATENÇÃO — MODELO PROVISÓRIO: ver a ressalva de status no topo de
% "relatorio-metalinguistico.tex" e em "../README.md". A TAL ainda não é
% modalidade de prática profissional confirmada para o PPC 2012.
%
% Este arquivo é o irmão de "relatorio-metalinguistico.tex": mesmo modelo,
% mesma estrutura de seções, mas com uma tutoria fictícia preenchida do
% início ao fim, para servir de referência visual de como o relatório
% pronto deve ficar, SE a modalidade vier a ser aceita.
%
% ATENÇÃO: dados, números e nomes deste exemplo são inventados para fins
% didáticos. Não citar nem reutilizar como se fossem reais.
%
% Compila com pdflatex (rode duas vezes por causa do sumário/referências).
% ============================================================================

\documentclass[article,12pt,a4paper]{abntex2}

\usepackage[T1]{fontenc}
\usepackage[utf8]{inputenc}
\usepackage[brazil]{babel}
\usepackage{times}
\usepackage{indentfirst}
\usepackage{graphicx}
\usepackage{microtype}

\renewcommand{\maketitlehookb}{}%% sem título em língua estrangeira (sem abstract)
\titulo{Tutoria de Aprendizagem e Laboratório em Lógica de Programação: um Relato de Experiência}
\autor{Larissa Beatriz Nunes Cavalcante \and Prof.\ Eduardo Matias de Albuquerque\footnote{Professor orientador e supervisor da TAL. Docente da área de Informática, IFRN Campus Ceará-Mirim.}}
\local{Ceará-Mirim -- RN}
\data{2026}
\instituicao{Instituto Federal de Educação, Ciência e Tecnologia do Rio Grande do Norte -- IFRN \\ Campus Ceará-Mirim}

\begin{document}
\maketitle

\begin{resumo}
A disciplina de Fundamentos de Lógica e Algoritmos apresenta,
historicamente, um dos maiores índices de dependência no primeiro ano do
curso Técnico Integrado em Informática. Este relatório tem como objetivo
descrever a tutoria de aprendizagem e laboratório (TAL) realizada na
disciplina, voltada a alunos do primeiro ano com dificuldade nos conteúdos
iniciais de lógica de programação. A tutoria consistiu em sessões
semanais de reforço, em formato de plantão de dúvidas e resolução guiada
de exercícios, ao longo de um semestre letivo. Participaram das sessões,
ao menos uma vez, 24 dos 68 alunos matriculados na disciplina, com um
grupo de 11 alunos frequentando regularmente. Entre os alunos que
frequentaram regularmente a tutoria, a taxa de aprovação na disciplina foi
de 91\%, contra 64\% entre os alunos que não participaram. Conclui-se que a
tutoria contribuiu para a redução da dificuldade inicial dos alunos com a
disciplina, e proporcionou à tutora experiência prática de ensino e
reforço de seu próprio domínio do conteúdo.

\vspace{\onelineskip}
\noindent
\textbf{Palavras-chave}: Tutoria entre pares. Lógica de programação. Reforço escolar. Ensino técnico.
\end{resumo}

\textual

\section{Introdução}

A disciplina de Fundamentos de Lógica e Algoritmos, cursada no primeiro
ano do curso Técnico Integrado em Informática, introduz conceitos
abstratos — variáveis, estruturas de repetição, estruturas condicionais —
que costumam representar a primeira barreira formal de programação para
estudantes recém-chegados ao ensino técnico. Nos últimos períodos letivos,
a disciplina apresentou índice de reprovação superior a 30\%, segundo
registros do professor titular.

Essa dificuldade recorrente motivou a oferta de tutoria de aprendizagem e
laboratório (TAL) na disciplina, com o objetivo de dar suporte adicional
aos alunos com maior dificuldade, fora do horário regular de aula. A
relação com a área técnica do curso é direta: a tutora já havia cursado e
sido aprovada com bom desempenho na disciplina, e a tutoria reforça, para
ela própria, o domínio desse conteúdo de base para as disciplinas
seguintes de programação.

O objetivo geral desta tutoria foi reduzir a dificuldade inicial dos
alunos do primeiro ano com os conteúdos de lógica de programação. Os
objetivos específicos foram: (a) oferecer sessões semanais de apoio fora
do horário de aula; (b) atender, prioritariamente, os alunos indicados
pelo professor da disciplina como tendo maior dificuldade; e (c) medir o
efeito da tutoria sobre a taxa de aprovação dos participantes.

Este relatório está organizado da seguinte forma: a seção 2 apresenta os
conceitos da disciplina trabalhados na tutoria; a seção 3 descreve a
metodologia das sessões; a seção 4 apresenta os resultados; e a seção 5
traz as considerações finais.

\section{Fundamentação}

Os conteúdos trabalhados na tutoria foram os mesmos do início da ementa da
disciplina: variáveis e tipos de dados, operadores, estruturas
condicionais (\texttt{se}/\texttt{senão}) e estruturas de repetição
(\texttt{enquanto}, \texttt{para}), representados em pseudocódigo e
fluxogramas, conforme o material do professor titular da disciplina.

A tutoria entre pares é uma prática em que um estudante que já domina um
conteúdo apoia colegas com dificuldade nele, em formato mais próximo e
informal do que a aula regular, o que costuma reduzir a barreira de pedir
ajuda e permite um ritmo adaptado a cada aluno.

Os conceitos revisados nas sessões são pré-requisito direto para as
disciplinas de Programação Estruturada e Orientada a Objetos e Programação
com Acesso a Banco de Dados, do núcleo tecnológico do curso, o que situa a
tutoria como reforço de uma base necessária à formação técnica.

\section{Metodologia e Descrição das Atividades}

A tutoria foi planejada em sessões semanais de duas horas, às
quartas-feiras no turno oposto às aulas, no laboratório de informática do
câmpus, ao longo do segundo semestre letivo de 2026, conforme o plano de
atividades deferido pelo professor orientador.

O público-alvo prioritário eram os alunos do primeiro ano indicados pelo
professor titular da disciplina como tendo maior dificuldade nas primeiras
avaliações, mas a participação era aberta a qualquer aluno matriculado.
Ao todo, 24 dos 68 alunos matriculados frequentaram ao menos uma sessão,
com um grupo de 11 alunos frequentando regularmente (mais de 8 das 14
sessões realizadas).

Cada sessão seguia o formato de plantão de dúvidas: a tutora revisava
brevemente o conteúdo da semana e, em seguida, os alunos resolviam
exercícios da lista da disciplina com apoio individual da tutora, que
circulava entre os grupos esclarecendo dúvidas. Foram usadas as listas de
exercícios do próprio professor titular e um conjunto adicional de
exercícios de fixação elaborado pela tutora com apoio do orientador.

\section{Resultados e Impacto}

A Tabela \ref{tab:resultados-tal} compara o desempenho na disciplina entre
os alunos que frequentaram regularmente a tutoria e os que não
participaram.

\begin{table}[!ht]
  \centering
  \caption{Taxa de aprovação na disciplina, por participação na tutoria}
  \label{tab:resultados-tal}
  \begin{tabular}{lcc}
    \hline
    Grupo & Nº de alunos & Taxa de aprovação \\
    \hline
    Frequentaram regularmente a tutoria & 11 & 91\% \\
    Não participaram da tutoria         & 44 & 64\% \\
    \hline
  \end{tabular}
  \\[0.3em]
  {\footnotesize Fonte: elaborado pela autora (2026), a partir das listas de presença da tutoria e do diário da disciplina, com autorização do professor titular.}
\end{table}

Ao todo, foram realizadas 14 sessões de tutoria, com média de 6 alunos por
encontro. Entre os 11 alunos que frequentaram regularmente, 10 foram
aprovados na disciplina, uma taxa bem acima da observada entre os alunos
que não participaram das sessões. Em conversa informal ao final do
semestre, os alunos frequentes relataram que o formato de grupo pequeno e
o ritmo mais lento das sessões ajudaram a entender conceitos que não
haviam ficado claros na aula regular.

Esses resultados respondem ao objetivo geral da tutoria, de reduzir a
dificuldade inicial dos alunos com a disciplina, e confirmam a relevância
de oferecer esse tipo de apoio já no primeiro ano do curso.

\section{Considerações Finais}

Esta tutoria atendeu ao objetivo de oferecer suporte adicional aos alunos
com maior dificuldade na disciplina de Fundamentos de Lógica e Algoritmos,
com taxa de aprovação de 91\% entre os participantes regulares, contra 64\%
entre os demais alunos matriculados.

Para a tutora, a experiência reforçou seu próprio domínio dos conteúdos de
lógica de programação e desenvolveu habilidades de explicação e
organização didática não trabalhadas nas disciplinas regulares do curso —
o que caracteriza a articulação entre teoria e prática também do lado de
quem tutora, e não só de quem é tutorado.

Como limitação, a participação caiu nas últimas semanas do semestre,
coincidindo com o período de avaliações de outras disciplinas, o que
sugere ajustar o calendário da tutoria em ofertas futuras. Como
continuidade, recomenda-se oferecer a tutoria já nas primeiras semanas do
semestre letivo seguinte, antes da primeira avaliação da disciplina, e não
apenas depois de identificadas as dificuldades.

\postextual

\section*{Referências}

\newcommand{\referencia}[1]{\par\noindent\hangindent=1cm #1\par\medskip}

\referencia{ALBUQUERQUE, Eduardo M. \textit{Material didático de
Fundamentos de Lógica e Algoritmos}. Ceará-Mirim: IFRN, 2026.}


\end{document}
